\section{Implementation}
\label{sec:impl}

\begin{table}
    \centering
    \caption{DNN benchmarks used in our evaluation.}
    \vspace{\captionvspace}
    \footnotesize
    \label{tab:benchmarks}
    \begin{tabular}{l|l|l}
    \toprule
    {\bf Name} & {\bf Description} & {\bf Base Architecture} \\
    \midrule
    GQA & Group-query attention & \llama-3-70B~\cite{grattafiori2024llama3herdmodels}\\
    QKNorm & QK normalization with attention & Chameleon-7B~\cite{chameleonteam2024chameleon} \\
    RMSNorm & RMS normalization with linear & \llama-2-7B~\cite{llama2}\\
    LoRA & Low-rank adaptation & GPT-3-7B-LoRA~\cite{llama-7b-lora} \\
    GatedMLP & Gated multi-layer perceptron & Falcon-7B~\cite{falcon40b}\\
    nTrans & Normalized Tarnsformer & nGPT-1B~\cite{loshchilov2024ngpt}\\

    \bottomrule
    \end{tabular}
    \vspace{\captionvspace}
    \vspace{\captionvspace}
\end{table}

\sys is implemented in 30K lines of code in C++, CUDA, and Python. Kernel operators are implemented with the cuDNN and cuBLAS libraries~\cite{cudnn, cublas}, and block and thread operators are implemented using cuTLASS~\cite{cutlass} and CUDA PTX.
%
For each input tensor program, \sys automatically generates and verifies potential \graphs. 
For each verified \graph, \sys produces CUDA source code for all custom kernels of the \graph and compiles the code into binary using the CUDA compiler.
%
This approach enables just-in-time (JIT) compilation and deployment for general tensor programs, and the generated kernels can be directly integrated into a PyTorch program with a few lines of code changes. 
%
\sys's SMT and ILP solvers are implemented using Z3 4.12.6~\cite{z3}. %\oded{I moved this here}

Our implementation supports the operators listed in \Cref{tab:expression}.
%\oded{Is this list needed? Can't we just say: ``Our implementation supports the operators listed in \Cref{tab:expression}''. If we still want to list everything in the text, I still think we should refer to the table.}
%\oded{What about repeat? It appears in Fig. 3}
%
\sys can be extended to include new operators, such as variants of convolution or matrix multiplication, at the kernel, block, and/or thread levels.
%For each new operator, \sys requires (1) an implementation that computes the operator's output over finite fields, and (2) a floating point implementation of the operator at the kernel, block, and/or thread levels. (1) is used by the probabilistic equivalence verifier for examining \graph equivalence, and (2) is used for generating optimized tensor programs.
%Moreover, to allow the \graph generator to discover \graphs containing the new operator, \sys needs an extension to the equivalence axioms $\Aeq$ regarding the properties for the new operator.
%\MW{I explained how to adapt the generator for the new operator. Do you think we need more?}\ZJ{I think this is sufficient}
%\oded{I suggest the following rewrite to the last few sentences:
To support a new linear operator, \sys requires (1) a float-pointing implementation of the operator at the kernel, block, and/or thread levels, which is used by the \graph optimizer to generate CUDA kernels; (2) an implementation of the operator over modular arithmetic (see \S\ref{sec:verify}); and (3) an extension to the abstract expression axioms $\Aeq$ and $\Asub$ for the operator (see \S\ref{subsec:abstract_expression}).
%}
%

To utilize \Cref{thm:zero_test,thm:prob}, random tests should be performed with sufficiently large prime numbers $p$ and $q$ and iterated multiple times. Our current implementation uses the largest values of $p$ and $q$ whose product fits in 16-bit integers (i.e., $p=227, q=113$) to run these random tests on GPUs. We leverage \sys's GPU optimizations--such as keeping intermediate results in shared memory--to accelerate the search procedure.
We also perform a single random test without iterating it and compare all elements of the output tensors.
We note that this equivalence verification procedure does not introduce false negatives.
While it could, in theory, introduce false positives, we have not observed any in practice.
For these reasons, we consider this procedure sufficient for the search process and plan to add a final verification step that provides the theoretical guarantees only for the best \graph at the end of the optimization process.

\paragraph{Equivalence verification for non-\lax programs.} While \sys can generate \graphs for arbitrary tensor programs, the probabilistic equivalence verifier is limited to \lax programs and does not support certain DNN operators such as ReLU~\cite{relu}. 
As an alternative, we have developed a solver-based verifier for arbitrary tensor programs. The verifier relies on user-provided mathematical properties of individual operators (e.g., linearity, associativity, commutativity, and distributivity) defined in first-order logic and uses these properties to verify equivalence using an automated theorem prover. Compared to the probabilistic equivalence verifier, the solver-based verifier supports more general programs, while requiring additional manual effort to specify the properties of each new operator. A detailed discussion of the solver-based verifier is beyond the scope of this paper.

\if 0
One of our benchmarks, MLP, uses the {\tt ReLU} non-linear operator, which is not natively supported by \sys.
To apply \sys to MLP, we replace {\tt ReLU} by exponentiation, which is also a non-linear function and is not used by MLP. We then change exponentiation in the resulting optimized \graph back to {\tt ReLU} and verify equivalence by manual examination\footnote{Alternatively, we have developed a solver-based method to verify arbitrary non-\lax programs in first order logic, which is beyond the scope of this paper.}.
\fi